\subsection{Concepts}
\textbf{Colloid:} Glue-like, unable to pass through fine membranes. \\
\textbf{Colloidal Domain:} Dimensions from \SI{1}{\nm} to \SI{1}{\micro\m}. \\
\textbf{Colloidal Solution:} Intermediate between true solutions ($< \SI{1}{\nm}$) and suspensions ($> \SI{1000}{\nm}$). \\
\textbf{Phases:}
\begin{itemize}
    \item \textbf{Dispersed phase:} The particles.
    \item \textbf{Continuous phase:} The medium.
    \item \textbf{Interfacial phase:} Boundary layer with different properties from bulk.
\end{itemize}
\textbf{Monodisperse:} All particles are same size. \\
\textbf{Polydisperse:} Particles are of many different sizes.

\subsection{Stability}
\textbf{Thermodynamic Stability:} System is in lowest free energy state ($dF=0$). \\
\textbf{Kinetic Stability:} Metastable state with high energy barrier preventing coagulation. \\
\textbf{Coagulation:} Particles coming together, leading to phase splitting.

\subsection{Classification}
\textbf{Lyophilic (Solvent-loving):} Soluble macromolecules, thermodynamically stable (e.g., proteins, polymers). \\
\textbf{Lyophobic (Solvent-fearing):} Thermodynamically unstable, require stabilization (e.g., gold sol).

\subsubsection{Types of Colloidal Dispersions}
\begin{center}
\begin{tabular}{p{0.15\linewidth} p{0.25\linewidth} p{0.45\linewidth}}
\toprule
\textbf{Disp. Phase} & \textbf{Disp. Medium} & \textbf{Name (Examples)} \\
\midrule
Liquid & Gas & Liquid aerosol (Fog, mist, spray) \\
Solid & Gas & Solid aerosol (Smoke, dust) \\
Gas & Liquid & Foam (Whipped cream) \\
Liquid & Liquid & Emulsion (Milk, mayonnaise) \\
Solid & Liquid & Sol, suspension (Paint, ink, gold sol) \\
Gas & Solid & Solid foam (Insulating foam, pumice) \\
Liquid & Solid & Solid emulsion (Pearl, cheese, bitumen) \\
Solid & Solid & Solid suspension (Ruby glass, alloys) \\
\bottomrule
\end{tabular}
\end{center}

\subsection{Fundamental Forces}
\begin{enumerate}
    \item \textbf{Gravitational:} Sedimentation or creaming (density difference).
    \item \textbf{Viscous Drag:} Resistance to motion.
    \item \textbf{Kinetic Energy:} Brownian motion (dominates for small particles).
    \item \textbf{Van der Waals:} Attractive force.
    \item \textbf{Electrostatic:} Repulsive force (stabilizing).
\end{enumerate}

\subsection{Characterization}
\textbf{Average Diameter ($\\bar{d}$):}
\begin{equation}
    \\bar{d} = \frac{\\sum d_i \\phi_i}{\\sum \\phi_i}
\end{equation}
where $\\phi_i$ is the weighting factor (number, surface, volume, intensity).
\vspace{0.2cm}